% =====================================================================
%  sections/07_evaluation.tex  --  VII. EVALUATION
% =====================================================================
\section{EVALUATION}
\label{sec:evaluation}



% ---------------------------------------------------------------------
\subsection{Experimental Setup}
\label{sub:experimental-setup}

This section specifies the concrete model, runtime, and hosting configuration
used to develop and evaluate \projname{}.

\begin{itemize}
  \item \textbf{Language Model:} The natural-language-to-SQL engine
      runs on \href{https://huggingface.co/Snowflake/Arctic-Text2SQL-R1-7B}{Snowflake/Arctic-Text2SQL-R1-7B},
      a 7B-parameter model fine-tuned specifically for text-to-SQL
      generation. The model is served locally using vLLM, a
      high-throughput inference engine, rather than accessed through a
      third-party API. Inference was configured with a temperature of
      \texttt{0.2} and a maximum output length of \texttt{2048} tokens, to
      favor deterministic, syntactically precise SQL generation.
  \item \textbf{Python Environment:} The backend and AI pipeline run on
      Python 3.11, with dependency management and virtual environment
      handling via Poetry.
  \item \textbf{Hosting:} The system is deployed on a remote Ubuntu server
      equipped with a GPU, self-managed and accessed like a local
      machine rather than through a managed cloud AI service, allowing
      the LLM to run entirely within our own infrastructure.
\end{itemize}

% ---------------------------------------------------------------------
\subsection{Project Strength}
\label{sub:project-strength}

Our team developed \projname{} with a focus on usability, performance,
and intelligent automation. Its key strengths include:

\begin{itemize}
  \item \textbf{AI Chat Functionality:} Provides real-time chat
        functionality to enable users to ask questions using natural
        language to gain insights. This allows non-technical business users
        to retrieve complex database records instantly without learning SQL.
  \item \textbf{Data Privacy and Security:} Ensures that all financial
        data is kept internal and protected from unauthorized access. This
        guarantees that sensitive corporate financial metrics remain secure
        and fully compliant with strict data protection standards.
  \item \textbf{Forecasting Capability:} Provides accurate and
        confidence metric predictions of future financial trends to
        support planning and decision-making. This empowers managers to proactively
        plan operational budgets and allocate resources based on data-driven projections.
  \item \textbf{Customizable Report Generation:} Allows users to quickly
        create and customize reports, delivering clear insights for
        analysis and decision-making, even without technical skills. This drastically
        reduces manual reporting overhead and accelerates recurring executive reviews.
\end{itemize}

% ---------------------------------------------------------------------
\subsection{Project Weakness}
\label{sub:project-weakness}

Despite its advantages, \projname{} still has some limitations:

\begin{itemize}
  \item \textbf{Handling Complex Queries:} May struggle with generating
        accurate results for highly complex or ambiguous natural
        language queries. Multi-nested subqueries or vague business terms can cause the Text-to-SQL engine to produce incorrect SQL joins or misinterpret user intent.
  \item \textbf{Limited Query Scope:} \projname{} only supports
        questions related to the financial data stored within the
        system. The system cannot synthesize external market context or answer general domain questions outside the indexed relational database tables.
\end{itemize}